\chapter{Logistic Regression}

\section{Introduction and Core Motivation}

Logistic regression is a fundamental supervised learning algorithm utilized for classification problems. Unlike regression models that predict a continuous numerical quantity, logistic regression models the conditional probability that a given observation belongs to a particular class.

\subsection{Why Not Standard Linear Regression?}
In ordinary linear regression, the relationship between features and the target is modeled as:
\begin{equation}
Y = \beta_0 + \beta_1 X_1 + \dots + \beta_p X_p + \varepsilon
\end{equation}
Under this linear formulation, the predicted value $\hat{Y}$ can take any real value across the entire unbounded interval:
\begin{equation}
\hat{Y} \in (-\infty, +\infty)
\end{equation}
However, probabilities are mathematically constrained to lie strictly within the unit interval:
\begin{equation}
P(Y = 1 \mid X) \in [0, 1]
\end{equation}
If ordinary linear regression were used for classification, predictions could yield nonsensical probabilities less than $0$ or greater than $1$. Consequently, a transformation is needed that ``squashes'' the unbounded linear output $(-\infty, +\infty)$ into the valid probability range $[0, 1]$. In logistic regression, this transformation is provided by the \textbf{sigmoid (logistic) function}.

\section{The Sigmoid (Logistic) Function and Probability Modeling}

Logistic regression first computes a linear combination of the input features:
\begin{equation}
Z = \beta_0 + \beta_1 X_1 + \dots + \beta_p X_p
\end{equation}
The real-valued score $Z$ is then passed through the sigmoid (logistic) activation function to yield the predicted probability of the positive class ($Y=1$):
\begin{equation}
P(Y=1) = \frac{1}{1 + e^{-Z}}
\end{equation}

\subsection{Algebraic Equivalence of the Sigmoid Function}
Multiplying both the numerator and the denominator by $e^Z$ (since $\frac{e^Z}{e^Z} = 1$) provides the alternate standard representation:
\begin{align}
P(Y=1) &= \frac{1 \cdot e^Z}{(1 + e^{-Z}) \cdot e^Z} \\
&= \frac{e^Z}{e^Z + e^{-Z} \cdot e^Z} \\
&= \frac{e^Z}{e^Z + e^{Z - Z}} \\
&= \frac{e^Z}{e^Z + e^0} \\
&= \frac{e^Z}{e^Z + 1} = \frac{e^Z}{1 + e^Z}
\end{align}
Substituting the linear combination $Z = \beta_0 + \beta_1 X_1 + \dots + \beta_p X_p$ directly yields:
\begin{equation}
P(Y=1) = \frac{e^{\beta_0 + \beta_1 X_1 + \dots + \beta_p X_p}}{1 + e^{\beta_0 + \beta_1 X_1 + \dots + \beta_p X_p}}
\end{equation}

\subsection{Properties of the Sigmoid Curve and Decision Threshold}
The sigmoid curve possesses distinctive geometric characteristics:
\begin{itemize}
    \item \textbf{S-Shaped Geometry:} The curve exhibits a smooth, monotonic S-shape.
    \item \textbf{Asymptotic Bounds:} As $Z \to -\infty$, $P(Y=1) \to 0$. As $Z \to +\infty$, $P(Y=1) \to 1$.
    \item \textbf{Symmetric Midpoint:} At $Z = 0$, $e^0 = 1$, yielding exactly $P(Y=1) = \frac{1}{1+1} = 0.5$.
\end{itemize}

To convert continuous probability estimates into discrete class labels, a \textbf{decision threshold} is established (conventionally set to $0.5$):
\begin{equation}
\hat{Y} = 
\begin{cases}
1, & \text{if } P(Y=1) \ge 0.5 \quad (Z \ge 0) \\
0, & \text{if } P(Y=1) < 0.5 \quad (Z < 0)
\end{cases}
\end{equation}

\begin{figure}[htbp]
\centering
\begin{tikzpicture}[x=0.80cm, y=4.4cm]
    % Background subtle grid/regions
    \fill[darkteal!5] (0, 0.5) rectangle (4.3, 1.0);
    \fill[crimson!5] (-4.3, 0) rectangle (0, 0.5);

    % Asymptotes and Threshold lines
    \draw[dashed, gray!60, thick] (-4.3, 1.0) -- (4.3, 1.0) 
        node[above, font=\scriptsize\sffamily, text=gray!80!black, fill=white, inner sep=1.5pt] at (2.2, 1.0) {Asymptote $P = 1.0$};
    \draw[dashed, navyblue!60, thick] (-4.3, 0.5) -- (4.3, 0.5) 
        node[below=2pt, font=\scriptsize\sffamily, text=navyblue!80!black, fill=white, inner sep=1.5pt] at (2.6, 0.5) {Threshold $= 0.5$};

    % Axes
    \draw[->, >=stealth, thick, darkgray] (-4.7, 0) -- (4.8, 0) node[right, font=\small] {$Z$};
    \draw[->, >=stealth, thick, darkgray] (0, -0.06) -- (0, 1.18) node[above, font=\small] {$P(Y=1)$};

    % Axis ticks and values (with opaque white backgrounds preventing strikethrough)
    \foreach \x in {-4, -2, 2, 4} {
        \draw (\x, 0.02) -- (\x, -0.02) node[below=2pt, font=\scriptsize] {$\x$};
    }
    \draw (-0.08, 1.0) -- (0.08, 1.0);
    \node[anchor=east, font=\scriptsize, fill=white, inner xsep=2.5pt, inner ysep=1pt] at (-0.10, 1.0) {$1.0$};
    
    \draw (-0.08, 0.5) -- (0.08, 0.5);
    \node[anchor=east, font=\scriptsize, fill=white, inner xsep=2.5pt, inner ysep=1pt] at (-0.10, 0.5) {$0.5$};
    
    \draw (-0.08, 0) -- (0.08, 0);
    \node[anchor=north east, font=\scriptsize, fill=white, inner sep=1pt] at (-0.08, -0.02) {$0$};

    % Sigmoid curve
    \draw[navyblue, ultra thick, domain=-4.3:4.3, samples=120, smooth] 
        plot (\x, {1 / (1 + exp(-\x))});

    % Inflection / Decision boundary point
    \filldraw[crimson] (0, 0.5) circle (2.4pt);
    \node[above right=2pt, xshift=2pt, font=\scriptsize\bfseries\sffamily, text=crimson, fill=white, inner sep=1.5pt] at (0, 0.5) {$(0,\, 0.5)$};

    % Outer Curly Braces & Prediction Region Annotations
    % Upper right brace: Class 1 (spanning P in [0.5, 1.0])
    \draw [decorate, decoration={brace, amplitude=6pt}, darkteal, line width=1.1pt] 
        (4.4, 1.0) -- (4.4, 0.5) 
        node [midway, right=8pt, align=left] {
            \textbf{\sffamily\color{darkteal} Predict Class 1} \textsf{\scriptsize\color{darkteal}($\hat{Y} = 1$)}\\[2pt]
            \scriptsize\sffamily\color{darkgray} $P(Y=1) \ge 0.5 \iff Z \ge 0$
        };

    % Lower left brace: Class 0 (spanning P in [0.0, 0.5])
    \draw [decorate, decoration={brace, amplitude=6pt}, crimson, line width=1.1pt] 
        (-4.4, 0.0) -- (-4.4, 0.5) 
        node [midway, left=8pt, align=right] {
            \textbf{\sffamily\color{crimson} Predict Class 0} \textsf{\scriptsize\color{crimson}($\hat{Y} = 0$)}\\[2pt]
            \scriptsize\sffamily\color{darkgray} $P(Y=1) < 0.5 \iff Z < 0$
        };
\end{tikzpicture}
\caption{The Sigmoid activation function squashing unbounded linear values $Z$ into valid class probabilities $[0, 1]$ with a decision boundary at $Z = 0$ ($P = 0.5$).}
\label{fig:sigmoid_activation}
\end{figure}

\newpage
\section{Key Terms and Matrix Notation}

\begin{definition}{Core Variables of Logistic Regression}{lr_terms}
\begin{itemize}
    \item \textbf{Independent Variables ($X$):} The set of input features or predictor variables $X_1, X_2, \dots, X_p$.
    \item \textbf{Dependent Variable ($Y$):} The categorical target outcome, typically binary $Y \in \{0, 1\}$.
    \item \textbf{Logistic Function (Sigmoid Function):} The nonlinear squashing function $\sigma(Z) = \frac{1}{1 + e^{-Z}}$ mapping real scores to $(0, 1)$.
    \item \textbf{Coefficient ($\beta_j$):} The parameter estimated by the model for predictor $X_j$, indicating how strongly $X_j$ affects the dependent variable on the log-odds scale.
    \item \textbf{Intercept ($\beta_0$):} The constant baseline term representing the log-odds when all independent variables equal zero ($X_1 = X_2 = \dots = X_p = 0$).
\end{itemize}
\end{definition}

\subsection{Matrix Representation}
In compact linear algebra notation, the dataset of $n$ observations and $p$ features is represented as an $n \times p$ feature matrix $X$:
\begin{equation}
X = \begin{bmatrix}
x_{11} & x_{12} & \dots & x_{1p} \\
x_{21} & x_{22} & \dots & x_{2p} \\
\vdots & \vdots & \ddots & \vdots \\
x_{n1} & x_{n2} & \dots & x_{np}
\end{bmatrix}
\end{equation}
where:
\begin{itemize}
    \item $n$ is the total number of observations (samples).
    \item $p$ is the total number of features (predictors).
\end{itemize}
The observed class labels form a binary response vector $y \in \{0, 1\}^n$:
\begin{equation}
y = \begin{bmatrix} y_1 \\ y_2 \\ \vdots \\ y_n \end{bmatrix}, \quad y_i \in \{0, 1\}
\end{equation}

\section{Odds, Log-Odds (Logit), and Parameter Interpretation}

\subsection{Definition of Odds}
The \textbf{odds} of an event represent the ratio of the probability that the event occurs to the probability that the event does not occur:
\begin{equation}
\text{Odds} = \frac{P(Y = 1)}{1 - P(Y = 1)}
\end{equation}

\begin{example}{Calculating Probabilities and Odds}{odds_apples_oranges}
Consider a basket containing $3$ apples and $5$ oranges (total items $= 3 + 5 = 8$).

\begin{enumerate}
    \item \textbf{Probability of selecting an Orange:}
    \begin{equation}
    P(\text{Orange}) = \frac{5}{8} = 0.625
    \end{equation}

    \item \textbf{Odds of selecting an Orange:}
    \begin{equation}
    \text{Odds}(\text{Orange}) = \frac{\text{Number of Oranges}}{\text{Number of Not Oranges (Apples)}} = \frac{5}{3} \approx 1.667
    \end{equation}

    \item \textbf{Odds of selecting an Apple:}
    \begin{equation}
    \text{Odds}(\text{Apple}) = \frac{\text{Number of Apples}}{\text{Number of Not Apples (Oranges)}} = \frac{3}{5} = 0.600
    \end{equation}
\end{enumerate}
\end{example}

\subsection{Log-Odds (Logit Function)}
Taking the natural logarithm of the odds produces the \textbf{log-odds} (or \textbf{logit}):
\begin{equation}
\ln(\text{Odds}) = \ln\left(\frac{P(Y=1)}{1 - P(Y=1)}\right)
\end{equation}
Substituting $P(Y=1) = \frac{e^Z}{1 + e^Z}$ and $1 - P(Y=1) = \frac{1}{1 + e^Z}$:
\begin{equation}
\frac{P(Y=1)}{1 - P(Y=1)} = \frac{\frac{e^Z}{1+e^Z}}{\frac{1}{1+e^Z}} = e^Z
\end{equation}
Taking the natural logarithm:
\begin{equation}
\ln\left(\frac{P(Y=1)}{1 - P(Y=1)}\right) = \ln(e^Z) = Z = \beta_0 + \beta_1 X_1 + \dots + \beta_p X_p
\end{equation}
Thus, \textbf{logistic regression is a linear model for the log-odds}.

\subsection{Interpreting Coefficients and Odds Ratios}
For a one-unit increase in predictor $X_j$ (holding all other variables constant), the log-odds changes by $\beta_j$:
\begin{equation}
\Delta \ln(\text{Odds}) = \beta_j
\end{equation}
Exponentiating both sides yields the \textbf{Odds Ratio (OR)}:
\begin{equation}
\text{Odds Ratio} = e^{\beta_j}
\end{equation}

The directional impact of $\beta_j$ on the odds is summarized as follows:
\begin{itemize}
    \item $\beta_j > 0 \implies e^{\beta_j} > 1$: The odds of the event increase.
    \item $\beta_j < 0 \implies e^{\beta_j} < 1$: The odds of the event decrease.
    \item $\beta_j = 0 \implies e^{\beta_j} = 1$: Predictor $X_j$ has no effect on the odds.
\end{itemize}

\begin{example}{Odds Ratio Interpretation with $\beta_j = 0.7$}{coef_interp}
Suppose the estimated coefficient for a predictor $X_j$ is $\beta_j = 0.7$.
\begin{equation}
\text{Odds Ratio} = e^{\beta_j} = e^{0.7} \approx 2.0137 \approx 2
\end{equation}
\textbf{Interpretation:} A one-unit increase in $X_j$ roughly doubles the odds of the outcome event occurring ($Y=1$).
\end{example}

\section{Model Training via Maximum Likelihood Estimation (MLE)}

Logistic regression estimates its parameters $\beta_0, \beta_1, \dots, \beta_p$ using \textbf{Maximum Likelihood Estimation (MLE)}. 
\begin{itemize}
    \item \textbf{Objective of MLE:} Find the parameter vector $\beta$ that maximizes the joint likelihood of observing the given sample data.
    \item \textbf{Lack of Closed-Form Solution:} Unlike ordinary least squares (OLS) regression where $\hat{\beta} = (X^T X)^{-1} X^T y$, there is \textbf{no closed-form analytical solution} for the MLE of logistic regression parameters.
    \item \textbf{Numerical Optimization:} Parameter estimates are solved numerically using iterative optimization algorithms (e.g., Newton-Raphson, Iteratively Reweighted Least Squares, or Gradient Ascent).
\end{itemize}

\section{Assumptions of Logistic Regression}

Valid inference and stable performance in logistic regression rely on six foundational assumptions:
\begin{enumerate}
    \item \textbf{Independent Observations:} Sample observations must be independent of one another (no repeated measures or auto-correlation).
    \item \textbf{Binary Dependent Variable:} The response variable $Y$ is binary (for standard binomial logistic regression).
    \item \textbf{Linearity in Log-Odds:} Independent variables exhibit a strictly linear relationship with the log-odds $\ln\left(\frac{P}{1-P}\right)$.
    \item \textbf{Absence of Extreme Outliers:} The predictor space should not contain extreme leverage points or anomalous outliers.
    \item \textbf{Adequately Large Sample Size:} Because MLE relies on large-sample asymptotic properties, sufficient sample size is required to prevent biased parameter estimates.
    \item \textbf{Low Multicollinearity:} Predictor variables should exhibit little to no multicollinearity. High correlation among predictors inflates parameter standard errors.
\end{enumerate}

\section{Variants, Metrics, and Trade-Offs}

\subsection{Types of Logistic Regression}
Depending on the nature of the target variable, logistic regression is categorized into three types:
\begin{itemize}
    \item \textbf{Binomial Logistic Regression:} The target has exactly two classes (e.g., Yes / No, Default / Non-Default).
    \item \textbf{Multinomial Logistic Regression:} The target has three or more unordered nominal classes (e.g., Cat, Dog, Sheep).
    \item \textbf{Ordinal Logistic Regression:} The target has three or more ordered categories (e.g., Low, Medium, High).
\end{itemize}

\subsection{Model Evaluation Metrics}
Model performance for logistic regression is assessed using standard classification metrics:
\begin{itemize}
    \item \textbf{Accuracy:} Ratio of correctly classified instances to total instances.
    \item \textbf{Precision:} Ratio of true positives to all predicted positives ($\frac{\text{TP}}{\text{TP} + \text{FP}}$).
    \item \textbf{Recall (Sensitivity):} Ratio of true positives to all actual positives ($\frac{\text{TP}}{\text{TP} + \text{FN}}$).
    \item \textbf{F1-Score:} Harmonic mean of precision and recall ($2 \cdot \frac{\text{Precision} \cdot \text{Recall}}{\text{Precision} + \text{Recall}}$).
    \item \textbf{AUC-ROC:} Area Under the Receiver Operating Characteristic curve, measuring discrimination across all possible classification thresholds.
\end{itemize}

\subsection{Advantages and Disadvantages}

\begin{center}
\begin{tabular}{p{0.46\textwidth} p{0.48\textwidth}}
\toprule
\textbf{Advantages (Pros)} & \textbf{Disadvantages (Cons)} \\
\midrule
\begin{itemize}[leftmargin=*, nosep]
    \item Outputs well-calibrated probabilities, not just hard class labels.
    \item Highly interpretable parameters via odds ratios ($e^{\beta_j}$).
    \item Directly extends to multiclass classification problems.
\end{itemize}
&
\begin{itemize}[leftmargin=*, nosep]
    \item Interaction terms between features must be constructed manually.
    \item Interpretations are multiplicative on odds, not linear on probabilities.
    \item Complete separation prevents convergence (coefficients diverge to $\pm\infty$).
    \item Requires regularization (e.g., L1/L2 penalties) or informative priors to handle separation.
\end{itemize} \\
\bottomrule
\end{tabular}
\end{center}

\section{End-to-End Workflow and Generalized Linear Model (GLM) Classification}

\subsection{Five-Step Implementation Pipeline}
\begin{enumerate}
    \item \textbf{Prepare Data:} Encode categorical predictors (e.g., one-hot or dummy encoding) and scale continuous features if regularization is utilized.
    \item \textbf{Fit Model:} Estimate parameter vector $\beta$ using Maximum Likelihood Estimation (MLE) via numerical optimization.
    \item \textbf{Interpret Coefficients:} Exponentiate parameters to interpret predictor effects as odds ratios ($e^{\beta_j}$).
    \item \textbf{Predict and Threshold:} Compute predicted probabilities $P(Y=1) = \sigma(X\beta)$ and assign class labels based on a selected threshold (typically $0.5$).
    \item \textbf{Evaluate Model:} Measure performance using Accuracy, Precision, Recall, F1-Score, and AUC-ROC.
\end{enumerate}

\subsection{Membership in Generalized Linear Models (GLM)}
Logistic regression is a member of the broader \textbf{Generalized Linear Models (GLM)} family. In GLM theory, a model is defined by three components:
\begin{enumerate}
    \item \textbf{Random Component:} The dependent variable $Y$ follows a Bernoulli (or Binomial) distribution: $Y \sim \text{Bernoulli}(p)$.
    \item \textbf{Systematic Component:} A linear combination of predictors $Z = \beta_0 + \sum_{j=1}^p \beta_j X_j$.
    \item \textbf{Link Function:} The \textbf{logit link function} $g(p) = \ln\left(\frac{p}{1-p}\right) = Z$, which maps the expected probability parameter from $[0, 1]$ to the unbounded real line $(-\infty, +\infty)$.
\end{enumerate}
